% Experiment 1: Breadth First Search (BFS)
\subsection*{Experiment 1: Breadth First Search (BFS)}
\textbf{Script File:} \texttt{BFS.py} \\[4pt]
\textbf{Aim:} To implement Breadth-First Search (BFS) graph traversal using a queue data structure to visit vertices level by level. \\[4pt]
\textbf{Sample Input:} 4 edges: (A B), (A C), (B D), (C E); Starting node: A \\[6pt]
\begin{tcolorbox}[colback=black!5!white,colframe=black!75!white,title=Terminal Execution Output - Breadth First Search (BFS)]
\begin{verbatim}
---Build Your Graph---
How many edges (connection) does your graph has ?Enter each edge separated by a spacce(e.g., A B):
Edge 1Edge 2Edge 3Edge 4Enter the starting node for BFS: 
Your Graph Dictionary: {'A': ['B', 'C'], 'B': ['A', 'D'], 'C': ['A', 'E'], 'D': ['B'], 'E': ['C']}
Staring BFS Traversal...
Exploring node1: A
Exploring node1: B
Exploring node1: C
Exploring node1: D
Exploring node1: E
BFS traversal order: ['A', 'B', 'C', 'D', 'E']
\end{verbatim}
\end{tcolorbox}

